\documentclass[12pt]{article}
\usepackage[utf8]{inputenc}
\usepackage[T1]{fontenc}
\usepackage{hyperref}
\title{Lanas Lobu Language Resources for K-12 Curriculum Development}
\author{Compiled by Claude Sonnet 5.5 (claude-sonnet-5-5) for Andr\'as Kornai}
\date{}
\begin{document}
\maketitle
\section{Introduction}
Lanas Lobu (ISO 639-3 ruu) is an endangered Paitanic (Austronesian) language of Sabah, Malaysia. Almost nothing dedicated to it turned up in one search. Pages were not opened.

\section{Primary Educational Resources}
None found. No dictionary, grammar or school material was found.

\section{Linguistic and Reference Materials}
\begin{itemize}
\item Lobel, \emph{North Borneo Sourcebook} (University of Hawai\textquoteleft i Press, 2016). It covers Lobu Lanas among the Paitanic languages and includes wordlists. \url{https://www.oeaw.ac.at/resources/Record/9780824857820?lng=en}
\item Wikipedia article on the Kinabatangan languages, listing ruu. \url{https://en.wikipedia.org/wiki/Kinabatangan_language}
\item Glottolog entry for Lanas Lobu. \url{https://glottolog.org/resource/languoid/id/lana1236.bigmap.html}
\end{itemize}

\section{Religious and Community Resources}
None found.

\section{Digital and Audio Resources}
\begin{itemize}
\item Glosbe pairing ruu with Kinabalian (cbw). Its reliability for Lanas Lobu is unclear. \url{https://glosbe.com/ruu/cbw}
\end{itemize}

\section{Curriculum Development Context}
Teaching materials would have to be built from scratch, starting from the wordlists in Lobel's sourcebook and fieldwork with speakers.
\end{document}